\documentclass[aspectratio=169,10pt]{beamer}
\usepackage[T1]{fontenc}
\usepackage{amsmath,amssymb}

\definecolor{ink}{HTML}{203550}
\definecolor{teal}{HTML}{008783}
\definecolor{codebg}{HTML}{F3F6F8}
\definecolor{codecomment}{HTML}{5B6B73}
\setbeamercolor{frametitle}{fg=ink}
\setbeamercolor{structure}{fg=teal}
\setbeamercolor{normal text}{fg=ink,bg=white}
\setbeamertemplate{navigation symbols}{}
\setbeamertemplate{footline}{\hfill\scriptsize QX26AgenticDelegation\quad\insertframenumber/\inserttotalframenumber\hspace{4mm}\vspace{2mm}}
\setbeamersize{text margin left=8mm,text margin right=8mm}

\newcommand{\src}[1]{\vfill{\tiny\color{gray}#1}}
\newenvironment{leancode}{%
  \begin{beamercolorbox}[sep=1.2ex,wd=\textwidth]{block body}%
  \color{ink}\ttfamily\scriptsize
}{\end{beamercolorbox}}
\title{Formalización parcial de la dinámica de conocimiento}
\subtitle{AI, Human Cognition and Knowledge Collapse}
\author{fgalarzac\quad |\quad EconCSLib}
\date{Acemoglu, Kong \& Ozdaglar (2026)\quad NBER WP 34910}

\begin{document}

\begin{frame}[plain]
  \titlepage
  \centering
  \small Carpeta de trabajo: \texttt{QX26AgenticDelegation}
  \medskip

  \scriptsize Formalización generada y auditada con el flujo de EconCSLib.\
  El resultado se reporta como parcial.
\end{frame}

\begin{frame}{1. Qué se encontró en el artículo}
  \begin{columns}[T]
    \column{.56\textwidth}
    \textbf{Inventario de fuente}
    \begin{itemize}
      \item 30 presentaciones nombradas.
      \item Assumptions 1--2.
      \item Observations 1--2 y Propositions 1--16.
      \item Lemmas 1--2, A-1--A-5 y B-1--B-3.
    \end{itemize}
    \column{.40\textwidth}
    \textbf{Resultado de la corrida}
    \begin{itemize}
      \item Una frontera formalizada.
      \item Recurrencia de Proposition 1.
      \item Existencia y unicidad probadas.
      \item Estado final: \textbf{parcial}.
    \end{itemize}
  \end{columns}
  \medskip
  \textbf{Fuente fijada:} NBER Working Paper 34910, febrero de 2026.
  \src{SOURCE\_INVENTORY.md; FINAL\_VALIDATION\_REPORT.md.}
\end{frame}

\begin{frame}{2. La frontera matemática formalizada}
  En equilibrio simétrico, el esfuerzo agregado es $E_t=I e_t$. La precisión
  heredada sigue la recurrencia mostrada en Proposition 1:
  \[
    X_{t+1}
      =\left[\Sigma^2+
        \frac{1}{X_t+\lambda_G I e_t}\right]^{-1}.
  \]
  EconCSLib representa el estado como el par $(X_t,e_t)$ y define
  \[
    (X_{t+1},e_{t+1})
      =\operatorname{dynamicTransition}(X_t,e_t),
      \qquad e_{t+1}=e(X_{t+1},\tau_A).
  \]
  \begin{block}{Alcance exacto}
    Una vez fijadas la condición inicial y la función de esfuerzo no negativa,
    existe una única trayectoria que satisface la recurrencia.
  \end{block}
  \src{MainTheorems.lean; paper Proposition 1, pp. 17--18.}
\end{frame}

\begin{frame}{3. Formalización Lean: el Spec de Proposition 1}
\begin{leancode}
def paper\_proposition1\_dynamic\_recurrenceSpec : Prop :=\\
\quad $\forall$ (X1 SigmaSq lambdaG aggregation tauA : $\mathbb R$)\\
\qquad (effort : $\mathbb R\to\mathbb R\to\mathbb R$),\\
\quad $0\le X1\to$\\
\quad ($\forall$ X, $0\le$ effort X tauA) $\to$\\
\quad $\exists!$ path : $\mathbb N\to\mathbb R\times\mathbb R$,\\
\qquad IsDynamicRecurrence X1 SigmaSq lambdaG aggregation tauA effort path
\end{leancode}
  \begin{itemize}
    \item El \texttt{Spec} es transparente y constituye la superficie revisable.
    \item \texttt{existsUnique} exige una sola trayectoria para los parámetros dados.
    \item La hipótesis visible exige esfuerzo no negativo.
  \end{itemize}
  \src{PaperInterface.lean: paper\_proposition1\_dynamic\_recurrenceSpec.}
\end{frame}

\begin{frame}{4. Endpoint probado y estructura de la prueba}
\begin{leancode}
theorem paper\_proposition1\_dynamic\_recurrence :\\
\quad paper\_proposition1\_dynamic\_recurrenceSpec := by\\
\quad unfold paper\_proposition1\_dynamic\_recurrenceSpec\\
\quad intro X1 SigmaSq lambdaG aggregation tauA\\
\qquad effort \_ effort\_nonnegative\\
\quad exact dynamicRecurrence\_existsUnique X1 SigmaSq\\
\qquad lambdaG aggregation tauA effort effort\_nonnegative
\end{leancode}
  \begin{center}
    \small
    \texttt{dynamicTransition} $\Longrightarrow$
    \texttt{IsDynamicRecurrence}\par\smallskip
    $\Longrightarrow$ \texttt{dynamicRecurrence\_existsUnique}
    $\Longrightarrow$ \texttt{ProofInterface}
  \end{center}
  La existencia se construye con recursión sobre $\mathbb N$. La unicidad se
  prueba por inducción: dos trayectorias con el mismo estado inicial y la misma
  transición coinciden en cada periodo.
  \src{MainTheorems.lean; ProofInterface.lean.}
\end{frame}

\begin{frame}{5. Validación mecánica}
  \begin{columns}[T]
    \column{.48\textwidth}
    \begin{block}{Chequeo del paper}
      \centering
      \texttt{paper\_contribution.py check}\par\smallskip
      \Large\textbf{Exit code 0}\par\smallskip
      \normalsize 755 trabajos compilados
    \end{block}
    \column{.48\textwidth}
    \begin{block}{Endpoint exacto}
      \centering
      \texttt{lake build ProofInterface}\par\smallskip
      \Large\textbf{Exit code 0}\par\smallskip
      \normalsize 756 trabajos compilados
    \end{block}
  \end{columns}
  \medskip
  El endpoint tiene exactamente el tipo del \texttt{Spec}; no contiene
  \texttt{sorry}. Los archivos de auditoría, estado y trazabilidad permanecen
  dentro de la carpeta.
  \src{FAST\_CHECK\_RESULT.txt; PROOF\_BUILD\_RESULT.txt; audit/.}
\end{frame}

\begin{frame}{6. Resultado honesto y trabajo pendiente}
  \textbf{Probado:}
  \begin{itemize}
    \item existencia y unicidad de la trayectoria de la recurrencia determinista;
    \item no negatividad del componente de esfuerzo, dada la hipótesis sobre
      la función de mejor respuesta.
  \end{itemize}
  \textbf{No probado todavía:}
  \begin{itemize}
    \item existencia y unicidad del Perfect Bayesian Equilibrium simétrico;
    \item las otras 29 presentaciones nombradas;
    \item equivalencia semántica completa entre el paper y toda la clausura Lean;
    \item revisión independiente y cierre v11 completo.
  \end{itemize}
  \begin{alertblock}{Conclusión}
    La corrida certifica una pieza precisa de Proposition 1. No constituye una
    formalización completa del artículo ni del equilibrio económico.
  \end{alertblock}
  \src{FINAL\_VALIDATION\_REPORT.md; status.json.}
\end{frame}

\end{document}
